\documentclass[../../main.tex]{subfiles}

\begin{document}

\section{赋范线性空间}\label{section:ch07-赋范线性空间}

本章始终用 $E$ 表示实数的可测集. 定义 $\mathcal{F}$ 为在 $E$ 上可测且 a.e. 有限的所有扩充实值函数所组成的族. 定义 $\mathcal{F}$ 中的两个函数 $f$ 和 $g$ 为等价的, 记为 $f\cong g$, 若
\begin{align*}
\text{对几乎所有 }x\in E,\quad f(x)=g(x).
\end{align*}
这是一个等价关系, 即它是自反, 对称与传递的. 因此它诱导出 $\mathcal{F}$ 的由等价类的不交族组成的划分, 我们把它记为 $\mathcal{F}/\cong$. $\mathcal{F}/\cong$ 上存在自然的线性结构: 给定 $\mathcal{F}$ 中的两个函数 $f$ 和 $g$, 对它们的等价类 $[f]$ 和 $[g]$ 以及实数 $\alpha$ 和 $\beta$, 定义线性组合 $\alpha[f]+\beta[g]$ 为 $E$ 中使得 $f$ 和 $g$ 都有限的点 $x$ 取值 $\alpha f(x)+\beta g(x)$ 的 $\mathcal{F}$ 中函数的等价类. 这些线性组合是正确定义的, 因为它们与等价类的代表的选取无关. 该线性空间的零元素是在 $E$ 上 a.e. 消失的函数的等价类.

线性空间的子集称为一个子空间, 若它关于线性组合封闭. 存在 $\mathcal{F}/\cong$ 的子空间的自然簇 $\{L^{p}(E)\}_{1\leqslant p\leqslant\infty}$. 对于 $1\leqslant p<\infty$, 我们定义 $L^{p}(E)$ 为使得
\begin{align*}
\int_{E}|f|^{p}<\infty
\end{align*}
的等价类 $[f]$ 的族. 这是正确定义的, 这是由于若 $f\cong g$, 则 $\int_{E}|f|^{p}=\int_{E}|g|^{p}$. 对任何两个数 $a$ 和 $b$,
\begin{align*}
|a+b|\leqslant|a|+|b|\leqslant 2\max\{|a|,|b|\},
\end{align*}
因此
\begin{align}
|a+b|^{p}\leqslant 2^{p}\{|a|^{p}+|b|^{p}\}.\label{align:ch07-三角不等式估计}
\end{align}
我们从该不等式以及积分的线性与单调性推出: 若 $[f]$ 和 $[g]$ 属于 $L^{p}(E)$, 则线性组合 $\alpha[f]+\beta[g]$ 也如此. 因此 $L^{p}(E)$ 是一个线性空间. 当然, $L^{1}(E)$ 由可积函数的等价类组成.

我们称函数 $f\in\mathcal{F}$ 为本性有界的, 若存在某个 $M\geqslant 0$, 称为 $f$ 的本性上界, 使得
\begin{align*}
\text{对几乎所有 }x\in E,\quad |f(x)|\leqslant M.
\end{align*}
我们定义 $L^{\infty}(E)$ 为使得 $f$ 本性有界的等价类 $[f]$ 所组成的族. 容易看出这是正确定义的且 $L^{\infty}(E)$ 也是 $\mathcal{F}/\cong$ 的一个线性子空间.

为了简单性与方便, 我们称 $\mathcal{F}/\cong$ 中的等价类为函数且记为 $f$ 而非 $[f]$. 因此写 $f=g$ 意味着 $f-g$ 在 $E$ 上 a.e. 消失. 这个简化迫使我们定义与 $L^{p}(E)$ 空间有关的概念时有必要检验一致性. 例如, 断言 $L^{p}(E)$ 中的序列 $\{f_{n}\}$ 在 $E$ 上 a.e. 逐点收敛于函数 $f\in L^{p}(E)$ 是有意义的, 这是由于若对所有 $n$, $g_{n}\cong f_{n}$ 且 $f\cong g$, 则由于零测度集的可数族的并也是零测度集, 序列 $\{g_{n}\}$ 也在 $E$ 上 a.e. 逐点收敛到 $g$. 说 $L^{p}[a,b]$ 中的函数 $f$ 是连续的, 意味着存在连续函数在 $[a,b]$ 上与 $f$ a.e. 一致. 由于零测度集的补集在 $\mathbf{R}$ 中稠密, 仅有一个这样的连续函数, 而考虑这个唯一的连续函数作为 $[f]$ 的代表常常是方便的.

在 19 世纪后期观察到一个或多个实变量的实值函数是经典分析的基础要素, 考虑定义域是函数的线性空间的实值函数也是有用的: 这样的函数称为泛函. 显然, 为了系统地研究如序列的收敛性 (以及维持收敛序列的线性性质) 这样富有成效的概念, 进而定义连续泛函的概念, 将绝对值的概念从实数拓广到一般的线性空间是有用的\footnote{我们后面将看到连续性也可以在关于映射的定义域与值域的度量结构或更一般的拓扑结构中讨论.}. 这些研究中出现的一个概念称为范数.

\begin{definition}[范数]\label{definition:ch07-范数}
令 $X$ 为线性空间. $X$ 上的实值泛函 $\|\cdot\|$ 称为范数, 若对 $X$ 中的每个 $f$ 和 $g$ 以及每个实数 $\alpha$,
\begin{align*}
\text{(三角不等式)}\quad\|f+g\|\leqslant\|f\|+\|g\|,
\end{align*}
\begin{align*}
\text{(正齐次性)}\quad\|\alpha f\|=|\alpha|\|f\|,
\end{align*}
\begin{align*}
\text{(非负性)}\quad\|f\|\geqslant 0\quad\text{且}\quad\|f\|=0\text{ 当且仅当 }f=0.
\end{align*}
\end{definition}

谈到赋范线性空间, 我们指的是加上范数的线性空间. 若 $X$ 是一个赋予范数 $\|\cdot\|$ 的线性空间, 我们称 $X$ 中的一个函数是单位函数, 若 $\|f\|=1$. 对任何 $f\in X$, $f\neq 0$, 函数 $f/\|f\|$ 是单位函数: 它是 $f$ 的数乘, 我们称它为 $f$ 的规范化.

\begin{example}[赋范线性空间 $L^{1}(E)$]\label{example:ch07-L1范数}
对于 $L^{1}(E)$ 中的函数 $f$, 定义
\begin{align*}
\|f\|_{1}=\int_{E}|f|.
\end{align*}
则 $\|\cdot\|_{1}$ 是 $L^{1}(E)$ 上的一个范数. 事实上, 对于 $f$, $g\in L^{1}(E)$, 由于 $f$ 和 $g$ 在 $E$ 上是 a.e. 有限的, 我们从实数的三角不等式推出,
\begin{align*}
\text{在 }E\text{ 上 a.e. }|f+g|\leqslant|f|+|g|.
\end{align*}
因此, 根据积分的单调性与线性,
\begin{align*}
\|f+g\|_{1}=\int_{E}|f+g|\leqslant\int_{E}[|f|+|g|]=\int_{E}|f|+\int_{E}|g|=\|f\|_{1}+\|g\|_{1}.
\end{align*}
显然, $\|\cdot\|_{1}$ 是正齐次的. 最后, 若 $f\in L^{1}(E)$ 且 $\|f\|_{1}=0$, 则在 $E$ 上 $f=0$ a.e.. 因此 $[f]$ 是线性空间 $L^{1}(E)\subseteq\mathcal{F}/\cong$ 的零元, 即 $f=0$.
\end{example}

\begin{example}[赋范线性空间 $L^{\infty}(E)$]\label{example:ch07-Linfty范数}
对于 $L^{\infty}(E)$ 中的函数 $f$, 定义 $\|f\|_{\infty}$ 为 $f$ 的本性上界的下确界. 我们称 $\|f\|_{\infty}$ 为 $f$ 的本性上确界且宣称 $\|\cdot\|_{\infty}$ 是 $L^{\infty}(E)$ 上的范数. 正性和正齐次性用与前一个例子用过的相同方法得出. 为验证三角不等式, 我们首先证明 $\|f\|_{\infty}$ 是 $f$ 在 $E$ 上的本性上界, 即
\begin{align}
\text{在 }E\text{ 上},\quad|f|\leqslant\|f\|_{\infty}\text{ a.e.}.\label{align:ch07-Linfty本性上界}
\end{align}
事实上, 对每个自然数 $n$, 存在 $E$ 的子集 $E_{n}$ 使得
\begin{align*}
\text{在 }E\sim E_{n}\text{ 上 }|f|\leqslant\|f\|_{\infty}+1/n\text{ 且 }m(E_{n})=0.
\end{align*}
因此, 若我们定义 $E_{\infty}=\bigcup_{n=1}^{\infty}E_{n}$,
\begin{align*}
\text{在 }E\sim E_{\infty}\text{ 上 }|f|\leqslant\|f\|_{\infty}\text{ 且 }m(E_{\infty})=0.
\end{align*}
因此, $f$ 的本性上确界是 $f$ 的最小的本性上界, 即\eqref{align:ch07-Linfty本性上界}成立. 现在对 $f$, $g\in L^{\infty}(E)$, 对几乎所有 $x\in E$,
\begin{align*}
|f(x)+g(x)|\leqslant|f(x)|+|g(x)|\leqslant\|f\|_{\infty}+\|g\|_{\infty}.
\end{align*}
因此, $\|f\|_{\infty}+\|g\|_{\infty}$ 是 $f+g$ 的本性上界且因此
\begin{align*}
\|f+g\|_{\infty}\leqslant\|f\|_{\infty}+\|g\|_{\infty}.
\end{align*}
\end{example}

\begin{example}[赋范线性空间 $\ell_{1}$ 和 $\ell_{\infty}$]\label{example:ch07-lp序列}
有一族序列的赋范线性空间, 它具有较简单的结构但与 $L^{p}(E)$ 空间有许多相似之处. 对 $1\leqslant p<\infty$, 定义 $\ell_{p}$ 为满足
\begin{align*}
\sum_{k=1}^{\infty}|a_{k}|^{p}<\infty
\end{align*}
的实数序列 $a=(a_{1},a_{2},\cdots)$ 族. 不等式\eqref{align:ch07-三角不等式估计}表明 $\ell^{p}$ 中的两个序列的和也属于 $\ell^{p}$ 且显然实数乘以 $\ell^{p}$ 中的序列也属于 $\ell^{p}$. 因此 $\ell^{p}$ 是一个线性空间. 我们定义 $\ell^{\infty}$ 为实有界序列的线性空间. 对于 $\ell^{1}$ 中的序列 $a=(a_{1},a_{2},\cdots)$, 定义
\begin{align*}
\|\{a_{k}\}\|_{1}=\sum_{k=1}^{\infty}|a_{k}|.
\end{align*}
这是 $\ell^{1}$ 上的范数. 对于 $\ell^{\infty}$ 中的序列, 定义
\begin{align*}
\|\{a_{k}\}\|_{\infty}=\sup_{1\leqslant k<\infty}|a_{k}|.
\end{align*}
容易看出 $\|\cdot\|_{\infty}$ 是 $\ell^{\infty}$ 上的范数.
\end{example}

\begin{example}\label{example:ch07-Cab范数}
令 $[a,b]$ 为闭有界区间. 则 $[a,b]$ 上的连续实值函数的线性空间记为 $C[a,b]$. 由于 $[a,b]$ 上的每个连续函数取到最大值, 对于 $f\in C[a,b]$, 我们定义
\begin{align*}
\|f\|_{\max}=\max_{x\in[a,b]}|f(x)|.
\end{align*}
我们把证明它定义了一个我们称为最大值范数的范数留作练习.
\end{example}

\begin{exercise}\label{exercise:ch07-01}
对 $C[a,b]$ 中的 $f$, 定义
\begin{align*}
\|f\|_{1}=\int_{a}^{b}|f|.
\end{align*}
证明这是 $C[a,b]$ 上的一个范数. 也证明不存在数 $c\geqslant 0$ 使得对所有 $C[a,b]$ 中的 $f$,
\begin{align*}
\|f\|_{\max}\leqslant c\|f\|_{1},
\end{align*}
但存在 $c\geqslant 0$ 使得对所有 $C[a,b]$ 中的 $f$,
\begin{align*}
\|f\|_{1}\leqslant c\|f\|_{\max}.
\end{align*}
\end{exercise}

\begin{exercise}\label{exercise:ch07-02}
令 $X$ 为定义在 $\mathbf{R}$ 上的实系数的所有多项式组成的簇. 证明这是线性空间. 对于多项式 $p$, 定义 $\|p\|$ 为 $p$ 的系数的绝对值的和. 这是一个范数吗?
\end{exercise}

\begin{exercise}\label{exercise:ch07-03}
对于 $L^{1}[a,b]$ 中的 $f$, 定义 $\|f\|=\int_{a}^{b}x^{2}|f(x)|\,dx$. 证明这是 $L^{1}[a,b]$ 上的范数.
\end{exercise}

\begin{exercise}\label{exercise:ch07-04}
对于 $L^{\infty}[a,b]$ 中的 $f$, 证明
\begin{align*}
\|f\|_{\infty}=\min\{M\mid m\{x\in[a,b]\mid|f(x)|>M\}=0\}.
\end{align*}
此外, 若 $f$ 在 $[a,b]$ 上是连续的, 则
\begin{align*}
\|f\|_{\infty}=\|f\|_{\max}.
\end{align*}
\end{exercise}

\begin{exercise}\label{exercise:ch07-05}
证明 $\ell^{\infty}$ 和 $\ell^{1}$ 是赋范线性空间.
\end{exercise}

\end{document}